

\title{Math 354: Cosets and Equivalence Relations}

\section*{Cosets}

In Proposition 4.3 we obtained the conditions
\begin{align*}
(\forall g \in G) &: a \sim b \implies ga \sim gb, \tag{$*$} \\
\text{and} \quad (\forall g \in G) &: a \sim b \implies ag \sim bg,
\end{align*}
which lead to a complete description of all compatible relations on a group $G$ (where compatible relations means that the equivalence relation respects the group operation).

In fact, each of these two conditions leads to a description of the relation satisfying it. Note that we have a group structure on $G/{\sim}$ only if both are satisfied.

Let's begin with a description of ($*$):

\subsection*{Proposition 4.4:}
Let $\sim$ be an equivalence relation on a group $G$ satisfying ($*$). Then
\begin{itemize}
    \item the equivalence class of $e_G$ is a subgroup $H$ of $G$; and
    \item $a \sim b \iff a^{-1}b \in H \iff aH = bH$.
\end{itemize}

\noindent\underline{Proof}. \\
Let $H \subseteq G$ be the equivalence class of the identity; $H \neq \phi$ as $e_G \in H$. For $a, b \in H$, we have $e_G \sim b$ and hence $b^{-1} \sim e_G$ (by ($*$)); hence $ab^{-1} \sim a$ (by ($*$)) and hence
$$ab^{-1} \sim a \sim e_G$$
by the transitivity of $\sim$ and since $a \in H$. This shows that $ab^{-1} \in H$ for all $a, b \in H$, showing that $H$ is a subgroup.

Next, assume $a, b \in G$ and $a \sim b$. Multiplying on the left by $a^{-1}$, ($*$) implies $e_G \sim a^{-1}b$, that is $a^{-1}b \in H$. Since $H$ is closed under the operation this implies $a^{-1}bH \subseteq H$, hence $bH \subseteq aH$; as $\sim$ is symmetric, the same reasoning gives $aH \subseteq bH$, and so $aH = bH$. Hence we have that
$$a \sim b \implies a^{-1}b \in H \implies aH = bH.$$

\vspace{1em}

Lastly, assume $aH = bH$. Then $a = ae_G \in bH$ and hence $a^{-1}b \in H$. By definition of $H$, this means $e_G \sim a^{-1}b$. Multiplying on the left by $a$ gives that us $a \sim b$, thus completing the proof.

\noindent\underline{Remark}: Proposition 4.4 shows that the equivalence classes of an equivalence relation satisfying ($*$) are all of the form
$$aH$$
for a fixed subgroup $H$, as $a$ ranges in $G$.

\subsection*{Definition 4.5:}
The \textbf{left-cosets} of a subgroup $H$ in a group $G$ are the sets $aH$, for $a \in G$. The \textbf{right-cosets} of $H$ are the sets $Ha$, $a \in G$.

\vspace{1em}
\noindent A `Converse' to Proposition 4.4 is true:

\subsection*{Proposition 4.6:}
If $H$ is any subgroup of a group $G$, the relation $\sim_L$ defined by
$$\forall(a, b \in G) : a \sim_L b \iff a^{-1}b \in H$$
is an equivalence relation satisfying ($*$).

\noindent\underline{Proof}: It is a straightforward exercise to see that $\sim_L$ is an equivalence relation. To see that the relation satisfies ($*$), note that
$$a \sim_L b \implies a^{-1}b \in H \implies a^{-1}(g^{-1}g)b \in H \implies (ga)^{-1}(gb) \in H \implies ga \sim_L gb$$
for all $g \in G$.  Poroposition 4.4 and 4.6 togeter shows

\subsection*{Proposition 4.7:}
There is a one-to-one correspondence between subgroups of $G$ and equivalence relations on $G$ satisfying ($*$); for the relation $\sim_L$ corresponding to a subgroup $H$, $G/{\sim_L}$ may be described as the set of left-cosets $aH$ of $H$.

\vspace{1em}

\noindent It is an exercise to produce the proofs and mirror statements giving a description of all equivalence relations satisfying ($**$). By doing this one should obtain:

\subsection*{Proposition 4.8:}
There is a one-to-one correspondence between subgroups of $G$ and equivalence relations on $G$ satisfying ($**$); for the relation $\sim_R$ corresponding to a subgroup $H$, $G/{\sim_R}$ may be described as the set of right cosets $Ha$ of $H$.

\vspace{1em}

\noindent The relation corresponding to $H$ defined in this other way is given by
$$a \sim_R b \iff ab^{-1} \in H \iff Ha = Hb.$$

\noindent\underline{Note}: The relations $\sim_L$ and $\sim_R$ corresponding to the same subgroup $H$ may not coincide. That is, left-cosets and right-cosets of a subgroup need not be equal. More generally
$$(\forall h \in H) : hH = Hh = H.$$
Also
$$(\forall a \in G) : a \in aH \cap Ha;$$

\vspace{1em}

\noindent So if $aH = Hb$, then in fact necessarily $aH = Ha$. This holds always if $G$ is \underline{commutative}, but it is not true in general.
\noindent\underline{Example 4.9}: Let $G = S_3$, and let $H$ be the subgroup consisting of the identity and the $1 \longleftrightarrow 2$ switch:
$$H = \left\{ \begin{pmatrix} 1 & 2 & 3 \\ 1 & 2 & 3 \end{pmatrix}, \begin{pmatrix} 1 & 2 & 3 \\ 2 & 1 & 3 \end{pmatrix} \right\}.$$
Then
\begin{align*}
\begin{pmatrix} 1 & 2 & 3 \\ 3 & 1 & 2 \end{pmatrix} H &= \left\{ \begin{pmatrix} 1 & 2 & 3 \\ 3 & 1 & 2 \end{pmatrix}\begin{pmatrix} 1 & 2 & 3 \\ 1 & 2 & 3 \end{pmatrix}, \begin{pmatrix} 1 & 2 & 3 \\ 3 & 1 & 2 \end{pmatrix}\begin{pmatrix} 1 & 2 & 3 \\ 2 & 1 & 3 \end{pmatrix} \right\} \\
&= \left\{ \begin{pmatrix} 1 & 2 & 3 \\ 3 & 1 & 2 \end{pmatrix}, \begin{pmatrix} 1 & 2 & 3 \\ 1 & 3 & 2 \end{pmatrix} \right\}
\end{align*}
while
\begin{align*}
H \begin{pmatrix} 1 & 2 & 3 \\ 3 & 1 & 2 \end{pmatrix} &= \left\{ \begin{pmatrix} 1 & 2 & 3 \\ 1 & 2 & 3 \end{pmatrix}\begin{pmatrix} 1 & 2 & 3 \\ 3 & 1 & 2 \end{pmatrix}, \begin{pmatrix} 1 & 2 & 3 \\ 2 & 1 & 3 \end{pmatrix}\begin{pmatrix} 1 & 2 & 3 \\ 3 & 1 & 2 \end{pmatrix} \right\} \\
&= \left\{ \begin{pmatrix} 1 & 2 & 3 \\ 3 & 1 & 2 \end{pmatrix}, \begin{pmatrix} 1 & 2 & 3 \\ 3 & 2 & 1 \end{pmatrix} \right\}
\end{align*}

\noindent This reflects the fact that the two conditions ($*$) and ($**$) are different.